% ===== PRÉAMBULE CANONIQUE — à coller VERBATIM en tête de CHAQUE .tex =====
\documentclass[11pt,a4paper]{article}
\usepackage[utf8]{inputenc}\usepackage[T1]{fontenc}\usepackage{lmodern}
\usepackage[french]{babel}
\usepackage{amsmath,amssymb,mathtools}\usepackage{icomma}
\usepackage{siunitx}\sisetup{locale=FR}  % \qty{}{}, \si{}, \ang{}
\usepackage{xcolor}\usepackage{colortbl}
\usepackage{tikz}\usepackage{pgfplots}\pgfplotsset{compat=1.18}
\usepackage[siunitx,europeanresistors]{circuitikz} % circuits (élec.)
\usepackage[most]{tcolorbox}\usepackage{enumitem}\usepackage{array,booktabs}
\usepackage[a4paper,margin=2.2cm]{geometry}
\usetikzlibrary{arrows.meta,calc,angles,quotes,positioning,patterns,%
  backgrounds,shapes.geometric,decorations.pathmorphing,decorations.markings,intersections}
\definecolor{primary}{RGB}{222,49,99}\definecolor{primarydark}{RGB}{150,22,55}
\definecolor{defcol}{RGB}{0,114,178}\definecolor{methcol}{RGB}{52,92,148}
\definecolor{excol}{RGB}{0,135,90}\definecolor{attncol}{RGB}{200,40,55}
\tcbset{boxbase/.style={enhanced,breakable,boxrule=0.9pt,arc=2.5pt,
  left=3mm,right=3mm,top=2.3mm,bottom=2.3mm,fonttitle=\bfseries,coltitle=white}}
% --- boîtes TUTORIEL (noms EXACTS, ne pas renommer) ---
\newtcolorbox{cle}[1][]{boxbase,colback=defcol!6,colframe=defcol,
  title={\textsc{À retenir}\ifx\relax#1\relax\relax\else\textmd{ : \itshape #1}\fi}}
\newtcolorbox{methode}[1][]{boxbase,colback=methcol!7,colframe=methcol,sharp corners,
  title={\textsc{Méthode}\ifx\relax#1\relax\relax\else\textmd{ --- \itshape #1}\fi}}
\newtcolorbox{exemplebox}[1][]{boxbase,colback=excol!5,colframe=excol,
  title={\textsc{Exemple}\ifx\relax#1\relax\relax\else\textmd{ : \itshape #1}\fi}}
\newtcolorbox{piege}[1][]{boxbase,colback=attncol!6,colframe=attncol,
  title={\textsc{Piège}\ifx\relax#1\relax\relax\else\textmd{ : \itshape #1}\fi}}
\newtcolorbox{remarque}[1][]{boxbase,colback=black!4,colframe=black!45,
  title={\textsc{Remarque}\ifx\relax#1\relax\relax\else\textmd{ : \itshape #1}\fi}}
% --- boîtes EXERCICE / CORRIGÉ (noms EXACTS, auto-comptées) ---
\newtcolorbox[auto counter]{exo}[1][]{boxbase,colback=primary!4,colframe=primary,
  title={\textsc{Exercice}~\thetcbcounter\ifx\relax#1\relax\relax\else\textmd{ --- \itshape #1}\fi}}
\newtcolorbox[auto counter]{corrige}[1][]{boxbase,colback=excol!4,colframe=excol,
  title={\textsc{Corrigé}~\thetcbcounter\ifx\relax#1\relax\relax\else\textmd{ --- \itshape #1}\fi}}
\newcommand{\vv}[1]{\vec{#1}}\newcommand{\ds}{\displaystyle}
\newcommand{\R}{\mathbb{R}}\newcommand{\N}{\mathbb{N}}\newcommand{\Z}{\mathbb{Z}}
\begin{document}
\sisetup{per-mode=symbol}

\begin{corrige}[vrai ou faux sur le microscope]
\begin{enumerate}[label=\alph*)]
  \item \textbf{Vrai} : image intermédiaire dans le plan focal objet de l'oculaire $\Rightarrow$
        rayons parallèles en sortie $\Rightarrow$ image finale à l'infini.
  \item \textbf{Vrai} : $G_{\text{oculaire}}\propto 1/f_2$, donc doubler $f_2$ divise le grossissement
        par deux.
  \item \textbf{Faux} : l'image de l'objectif est \emph{renversée} (et réelle, agrandie).
  \item \textbf{Faux} : l'oculaire est \emph{convergent} (il joue le rôle de loupe).
  \item \textbf{Vrai} : l'oculaire donne une image virtuelle, droite (par rapport à l'intermédiaire)
        et agrandie.
\end{enumerate}
Affirmations correctes : \textbf{a)}, \textbf{b)} et \textbf{e)}.
\end{corrige}

\begin{corrige}[la loupe brisée]
\textbf{Oui, les deux morceaux restent utilisables.} Chaque point de l'objet envoie de la lumière sur
\emph{toute} la surface de la lentille ; il suffit donc d'un fragment pour reconstituer l'image
\textbf{entière} de l'objet. En revanche, un morceau plus petit collecte \emph{moins de lumière} :
l'image est complète mais \textbf{moins lumineuse}.
\end{corrige}

\begin{corrige}[choisir le bon oculaire]
\begin{enumerate}[label=\alph*)]
  \item $G=|\gamma_1|\times G_{\text{oculaire}}=30\times 10=300$.
  \item $G_{\text{oculaire}}\propto 1/f_2$ : doubler $f_2$ le divise par deux, donc
        $G=300/2=150$.
  \item Parce que $G_{\text{oculaire}}\propto 1/f_2$ : plus la focale de l'oculaire est \emph{courte},
        plus son grossissement (et donc le grossissement total) est \emph{grand}.
\end{enumerate}
\end{corrige}

\begin{corrige}[photo ou projecteur ?]
\begin{enumerate}[label=\alph*)]
  \item $\dfrac{1}{v}=\dfrac{1}{5}-\dfrac{1}{200}=\dfrac{40-1}{200}=\dfrac{39}{200}\Rightarrow
        v\approx\qty{5.1}{\cm}$ ; $\gamma=-\dfrac{5,1}{200}\approx-0{,}026$. Image réelle, renversée,
        \textbf{très réduite}. Comme $u>2f$ : \textbf{appareil photo}.
  \item $\dfrac{1}{v}=\dfrac{1}{5}-\dfrac{1}{6}=\dfrac{6-5}{30}=\dfrac{1}{30}\Rightarrow v=\qty{30}{\cm}$ ;
        $\gamma=-\dfrac{30}{6}=-5$. Image réelle, renversée, \textbf{agrandie} $5\times$. Comme
        $f<u<2f$ : \textbf{projecteur}.
\end{enumerate}
\end{corrige}

\end{document}
